% Generated draft: requires mathematical and source review.
\documentclass[11pt]{article}
\usepackage[margin=25mm]{geometry}
\usepackage{fontspec}
\setmainfont{Latin Modern Roman}
\usepackage{amsmath,amssymb,mathtools,unicode-math}
\setmathfont{Latin Modern Math}
\usepackage{enumitem,needspace,xurl,hyperref}
\setlength{\parindent}{0pt}
\setlength{\parskip}{5pt}
\setlength{\emergencystretch}{3em}
\newcommand{\N}{\mathbb{N}}
\newcommand{\Z}{\mathbb{Z}}
\newcommand{\Q}{\mathbb{Q}}
\newcommand{\R}{\mathbb{R}}
\newcommand{\C}{\mathbb{C}}
\newcommand{\F}{\mathbb{F}}
\newcommand{\A}{\mathbb{A}}
\newcommand{\PP}{\mathbb{P}}
\newcommand{\E}{\mathbb{E}}
\newcommand{\eps}{\varepsilon}
\newcommand{\dd}{\,\mathrm{d}}
\begin{document}
\section*{UnsolvedMath ID: 20001931}
\textbf{Source problem code:} catalogue entry
\subsection{Definitions and mathematical statement}
Let \(S\) be a closed oriented surface of genus at least two. A regular smooth strongly convex Finsler metric is a function \(F:TS\to[0,\infty)\), positive off the zero section, smooth there, positively homogeneous of degree one, with positive-definite fundamental tensor
\[
 g_{ij}(x,y)=\frac12\frac{\partial^2F^2}{\partial y^i\partial y^j}(x,y).
\]
It is Riemannian precisely when this tensor is independent of \(y\), equivalently its Cartan tensor \(C_{ijk}=\tfrac12\partial g_{ij}/\partial y^k\) vanishes. Using the Chern horizontal covariant derivative, the Landsberg condition is
\[
 L_{ijk}=C_{ijk\mid s}y^s=0\quad\text{on }TS\setminus0.
\]
Question \(10.1.4\) asks whether such a surface admits a non-Riemannian Landsberg metric. The closed oriented regular strongly convex setting is the aid's explicit interpretation of the shorter source sentence; generalized singular Finsler structures are not included.

The source refers to [Bao2007] for the Landsberg condition and explains that it supplies a geodesic-flow integral, motivating its placement in the integrability section. It cites [Gom-Rug2013] for a negative answer under absence of conjugate points. The supplied aid combines these and further restrictions into the following necessary conditions for any positive example:
\begin{enumerate}
\item \(F\) is non-Berwald. Berwald means that its geodesic spray coefficients are quadratic in the tangent variable, or equivalently that its Berwald curvature vanishes. Thus an example would be a regular non-Berwald Landsberg metric, often called a regular unicorn.
\item \(F\) is geodesically complete, in both time directions, as a regular metric on a closed manifold; no reversibility assumption \(F(x,y)=F(x,-y)\) is made.
\item \(F\) has conjugate points, by the reported no-conjugate-points obstruction.
\item If \(K(p,v)\) is its surface flag curvature, there exist base points \(p_+,p_-\in S\) such that
\[
 K(p_+,v)>0\quad\text{for every }v\in T_{p_+}S\setminus0,
 \qquad
 K(p_-,v)<0\quad\text{for every }v\in T_{p_-}S\setminus0.
\]
The all-directions quantifiers are stronger than merely finding two tangent directions with opposite signs.
\end{enumerate}
These are reported obstructions for sufficiently smooth regular metrics, not an existence construction or a proof that the general question has a negative answer. The exact minimal differentiability and full theorem inputs for the basewise curvature assertion are not supplied. The extracted page header and the following Section \(10.2\) heading, ``Integrable geodesic flows with integrals of higher degree,'' are neighboring material rather than another question here.
\subsection{Short English statement}
The source asks for a regular non-Riemannian Landsberg metric on a higher-genus surface. The aid says any example must be non-Berwald, complete without reversibility, and have conjugate points. It also requires base points with flag curvature positive in every direction and negative in every direction, respectively; no example is supplied.
\subsection{Sources}
\begin{itemize}
\item Local UnsolvedMath registry, permanent ID \(20001931\), source code \texttt{catalogue entry}; canonical collection \texttt{aim-geometry-notes.json}, notes[268]. Source item \(10.1\), actual question \(10.1.4\).
\item Keith Burns and Vladimir S. Matveev, ``Open problems and questions about geodesics,'' AIM workshop ``Geodesics'': \url{https://aimath.org/pastworkshops/geodesicsproblems.pdf}. The supplied verification reports comparison with the PDF; the author/page header and Section \(10.2\) are extraction spillover. No PDF was opened in this session.
\item [Bao2007] for the Landsberg definition and [Gom-Rug2013] for the no-conjugate-points negative result. Full bibliographic entries and the additional curvature-sign theorem sources are not supplied. Original remarks and literature fields are empty; no proposer or year is supplied.
\item Supplied research attempt \(1\), \texttt{partial\_result}, run status ``open,'' rigor label ``proved,'' novelty confidence ``low.'' Registry status \texttt{partially\_solved}, classification \texttt{PARTIAL-PROGRESS}, updated 14 August 2026.
\end{itemize}
\end{document}
